% =====================================================================
% Chapter 9 -- Discussion
% =====================================================================
\documentclass[../preamble.tex]{subfiles}
\begin{document}

\chapter{Discussion}
\label{ch:discussion}

\section{What the experiments show}
\label{sec:discussion-main}

The central result is not simply an AUC of 0.9487. It is the gap between the
internal and external measurements. The foundation models reach AUC values
in the 0.99 range on the Figshare test split, but the same checkpoints score
0.45--0.59 on zero-shot PCOSgen evaluation. The confusion-matrix progression
shows why: the foundation model predicts PCOS for almost every target image.
This is a domain-shift failure, not an isolated bad architecture.

Fine-tuning on the PCOSgen train pool recovers the external signal. The best
single model reaches AUC 0.9420, and probability averaging the three most
complementary models reaches 0.9487. The ensemble reduces false positives
from 241 for DenseNet-121 alone to 164 while retaining sensitivity above
0.99. This is the strongest practical argument for using an ensemble: it
improves the operating profile without claiming that any single architecture
is universally superior.

\section{Why padding matters}
\label{sec:discussion-padding}

Letterbox padding creates a stable border pattern. When the border is
correlated with a class in the source dataset, a model can exploit it instead
of the ovarian morphology. The border-attention figure and the
padding-versus-no-pad AUC figure support this interpretation: removing the
padding improves cross-dataset performance by approximately 0.25 absolute at
the foundation stage. This is a shortcut-learning result, not a claim that
padding is always harmful for every ultrasound dataset.

\section{Why HPO did not replace the pre-HPO model}
\label{sec:discussion-hpo}

Optuna selected useful learning rates, dropout values, and batch sizes, but
the post-HPO retrain did not improve the external ensemble. The target
validation split has only 640 images and is itself heavily imbalanced. A
small validation set can make a hyperparameter configuration appear
preferable even when its ranking does not transfer to the held-out test.
The study therefore reports the pre-HPO ensemble as the recommended model
rather than choosing a configuration solely because it had the best internal
HPO score.

This result is not an argument against HPO. It is an argument for larger
validation cohorts, repeated cross-validation, or nested external validation
when HPO decisions are expected to affect deployment.

\section{Clinical interpretation}
\label{sec:discussion-clinical}

The default ensemble threshold produces sensitivity 0.9927 and specificity
0.8215. At the Youden-J threshold, sensitivity is 0.9909 and specificity is
0.8324. At the sensitivity-targeted threshold, sensitivity is fixed at 0.95
and specificity becomes 0.8553. These operating points are more useful than
an accuracy number because they expose the clinical trade-off between missed
PCOS cases and unnecessary referrals.

The intended clinical role is screening support, not autonomous diagnosis.
A screening system may reasonably favour sensitivity, but its threshold
must be selected with the prevalence and cost structure of the deployment
site. A model calibrated on PCOSgen cannot be assumed calibrated after a
hospital changes its equipment, patient population, or acquisition protocol.
The referral concept therefore requires monitoring, local threshold
validation, and prospective evaluation.

\section{Uncertainty and explainability together}
\label{sec:discussion-trust}

MC-Dropout entropy identifies cases in which the model's stochastic posterior
varies across passes. Grad-CAM identifies the region that contributed most to
the prediction. Their combination is more informative than either alone:

\begin{itemize}
  \item low entropy with a central ovarian activation is a plausible
        high-confidence case;
  \item high entropy with a diffuse or border-focused activation is a useful
        candidate for review; and
  \item low entropy with a border-focused activation exposes a confident
        shortcut, which aggregate calibration cannot detect.
\end{itemize}

These are qualitative interpretation rules, not validated clinical decision
criteria. Grad-CAM is sensitive to the selected target layer and provides a
coarse heatmap. The thesis does not claim that a heatmap proves clinical
causality. Objective perturbation scoring and clinician-reader studies are
appropriate next steps \cite{samek2017,tjoa2021}.

\section{Failure analysis}
\label{sec:discussion-failure}

The false-positive count remains 164 for the final ensemble at threshold 0.5.
This matters because PCOSgen is healthy-heavy in its held-out test set, so a
small false-positive rate can still produce a substantial false-positive
count. The false-negative count is four, which explains the high sensitivity;
however, four false negatives on one public cohort must not be interpreted as
a clinical false-negative rate.

The Grad-CAM panels show two broad failure patterns: attention that extends
beyond the central ovarian region, and uncertain images whose follicle
visibility is ambiguous even to visual inspection. The dataset does not
contain sufficient clinical metadata to distinguish true biological ambiguity
from acquisition artefact. The failure analysis therefore describes visual
patterns without assigning a clinical cause.

\section{Limitations}
\label{sec:discussion-limitations}

\begin{enumerate}
  \item Only one external public cohort is evaluated. A second independent
        hospital cohort is needed to assess population and device shift.
  \item The PCOSgen release does not document a verified patient-level split.
        The thesis therefore cannot guarantee patient independence in that
        external cohort.
  \item The HPO validation set is small and imbalanced. Repeated or nested
        cross-validation is needed for stable model selection.
  \item The wall-clock estimate is approximately 5--6 hours based on the
        recorded training files; HPO, calibration, XAI, and uncertainty were
        not separately time-stamped.
  \item The XAI analysis uses Grad-CAM only. LRP, SHAP, region perturbation,
        and clinician-reader validation were not executed.
  \item No demographic subgroup analysis is possible from the released
        metadata. Fairness across age, ethnicity, equipment, and clinical
        setting remains untested.
  \item The reported model is not a diagnosis system. Clinical deployment
        would require prospective validation, regulatory review, and a
        defined human-in-the-loop workflow.
\end{enumerate}

\section{Reproducibility and auditability}
\label{sec:discussion-repro}

The strongest reproducibility feature is the separation between design and
run. The thesis does not present EMA, SWA, k-fold CV, LRP, or SHAP as if they
were part of the final numbers. Their status is documented in the appendices,
and every headline metric points back to a JSON or CSV artefact. This makes
the limitations visible rather than hiding them behind a single performance
number.

\end{document}
